\subsubsection{Does the camera confirm the scaling?}

The foot lift above is measured by the Dot foot sensor (on the top of the foot). The camera toe marker confirms the scaling but with a
smaller slope (Table~\ref{tab:validityslopes}): +8.8\,mm (leading) and +7.5\,mm (trailing) per cm of obstacle, against +10.1 and
+11.1\,mm for the sensor. The sensor rises more than the toe, especially on the trailing foot, which crosses heel-up. Crossing by
crossing (Table~\ref{tab:validityagree}), the sensor rise agrees well with the camera toe rise for the leading foot (bias +8\,mm,
$r = 0.82$, within person 0.90) and less for the trailing foot (bias +48\,mm, $r = 0.62$); the provisional rigid-body toe estimate
removes the trailing bias but not the scatter. Together with the flat clearance slope (Table~\ref{tab:lme}), people raise the toe by
slightly less than the obstacle height and accept a marginally smaller clearance for higher obstacles (not significant).

\begin{table}[htbp]
\centering
\caption{Foot lift per cm of obstacle measured by the camera toe marker and by the Dot (mixed model, 95\,\% CI).}
\label{tab:validityslopes}
\begin{adjustbox}{max width=\textwidth}
{\footnotesize
\begin{tabular}{lrr}
\toprule
\textbf{foot lift measure} & \textbf{leading limb, slope per cm (95\% CI)} & \textbf{trailing limb, slope per cm (95\% CI)} \\
\midrule
camera toe rise & +8.8 mm (+7.0 to +10.6) & +7.5 mm (+5.4 to +9.6) \\
IMU foot-sensor rise & +10.1 mm (+9.4 to +10.8) & +11.1 mm (+9.7 to +12.6) \\
IMU toe rise (rigid-body, provisional) & +10.1 mm (+9.4 to +10.8) & +10.5 mm (+9.1 to +11.8) \\
\bottomrule
\end{tabular}
}
\end{adjustbox}

\end{table}

\begin{table}[htbp]
\centering
\caption{Agreement of the Dot foot lift with the camera toe rise, per crossing (Bland--Altman bias, SD and 95\,\% limits; correlation
overall and within participant).}
\label{tab:validityagree}
\begin{adjustbox}{max width=\textwidth}
{\scriptsize
\begin{tabular}{llrrrrrrr}
\toprule
\textbf{IMU vs camera toe rise} & \textbf{limb} & \textbf{crossings} & \textbf{bias (mm)} & \textbf{SD (mm)} & \textbf{95\% limits (mm)} & \textbf{median abs. difference (mm)} & \textbf{r} & \textbf{within-person r (median)} \\
\midrule
IMU foot-sensor rise & lead & 260 & +8 & 66 & $-$122 to +137 & 30 & 0.82 & 0.90 \\
IMU foot-sensor rise & trail & 259 & +48 & 108 & $-$163 to +260 & 77 & 0.62 & 0.70 \\
IMU toe rise (rigid-body, provisional) & lead & 260 & +5 & 66 & $-$124 to +134 & 29 & 0.82 & 0.90 \\
IMU toe rise (rigid-body, provisional) & trail & 259 & $-$3 & 109 & $-$216 to +210 & 46 & 0.60 & 0.69 \\
\bottomrule
\end{tabular}
}
\end{adjustbox}

\end{table}

\subsubsection{No anticipation in the approach}

None of the variables available before the leading foot leaves the floor changes with obstacle height (Table~\ref{tab:anticipation}):
the two approach strides, the approach speed and the leading-foot stance all have slopes indistinguishable from zero and explain at
most 1\,\% of the within-person variance; the camera trailing-foot placement before the obstacle explains 4\,\%. The four IMU
approach variables classify the height at chance level (forest 41.6\,\%, logistic regression 30.8\,\%, permutation $p = 0.75$).
The adaptation to the obstacle height happens inside the crossing step itself, starting at the leading toe-off.

\begin{table}[htbp]
\centering
\caption{Anticipation: change per cm of obstacle of the variables before the crossing step (mixed model) and their within-person
$\eta^2$.}
\label{tab:anticipation}
\begin{adjustbox}{max width=\textwidth}
{\footnotesize
\begin{tabular}{p{6.2cm}rrrr}
\toprule
\textbf{variable} & \textbf{slope per cm} & \textbf{95\% CI} & \textbf{p} & \textbf{$\eta^2$ within person} \\
\midrule
stride two before the crossing stride (leading foot, mm) & $-$1.38 & $-$4.55 to +1.78 & 0.346 & 0.01 \\
stride just before the crossing stride (leading foot, mm) & +0.519 & $-$1.32 to +2.36 & 0.574 & 0.01 \\
approach speed (m/s) & $-$0.000239 & $-$0.00193 to +0.00145 & 0.752 & 0.01 \\
leading-foot stance before its crossing swing (s) & +0.000492 & $-$0.000593 to +0.00158 & 0.351 & 0.01 \\
trailing toe before the obstacle at its toe-off (camera, mm) & +3.29 & $-$3.64 to +10.2 & 0.312 & 0.04 \\
\bottomrule
\end{tabular}
}
\end{adjustbox}

\end{table}
